% GPT-6 Astra Pro; effort not exposed; /root/replacement; 22 November 2026.
\section[HD-F1-R09: pure-character retracts and the balanced RM trace]{HD-F1-R09: one pure exponential summand obstructs the balanced RM trace}
\label{sec:followup-rm-pure-summand}

Let $(X,F,\theta)$ and $v_q$ be the genuine real-quadratic fourfold
and balanced character of HD-F1-R07. Consider any actual perfect
object with whole character
\[
 \operatorname{ch}(E)=N e^{-B}v_q,\qquad
 N\ne0,\quad q>0,\quad B\in\operatorname{NS}(X)_{\mathbf Q}.
 \tag{F1R9.1}
\]
An actual direct summand $Q$ of $E$ with
\[
                 \operatorname{ch}(Q)=r e^\ell,\qquad
 r\ne0,\quad\ell\in\operatorname{NS}(X)_{\mathbf Q}
 \tag{F1R9.2}
\]
forces a nonzero invariant degree-two full-trace kernel for every
coherent faithful finite point action and every lift of $E$.
Consequently no such lift is weakly $G$-semiregular. The summand need
not be preserved by the chosen linearization. Its complement is
arbitrary; $Q$ need not be a simple sheaf or bundle, and no negative-Ext
vanishing is assumed.

More generally, (F1R9.2) may be replaced by the character of an actual
object obtained from such a nonzero exponential character by a finite
composition of the exact Fourier--Mukai, isogeny, line-tensor,
translation, shift and dual operations in R08. This includes any
shift of a line bundle, a simple semihomogeneous bundle, and any
skyscraper sheaf, as a direct summand. The assertion concerns an
actual retract and makes no assertion about new nonsplit extensions.

\paragraph{1. A quotient obstruction for the exact degree-two kernel.}
Use the normalization shear (F1R7.5). On the normalized balanced
character,
\[
 U=U_1\oplus U_2,\quad\dim U_i=4,\qquad
 K=\ker c_{v_q}^2=\ker\omega_1\oplus\ker\omega_2,
 \tag{F1R9.3}
\]
where both $\omega_i:\bigwedge^2U_i\to\mathbf C$ are nondegenerate
alternating forms. Every linear map $\pi:U\to W$ satisfying
$\bigwedge^2\pi(K)=0$ has rank at most two.

Here is the exact argument. Put $\pi_i=\pi|_{U_i}$. The restriction
$\bigwedge^2\pi_i$ factors through the one-dimensional quotient
$\bigwedge^2U_i/\ker\omega_i$, so its rank is at most one.
If $\operatorname{rank}\pi_i\ge3$, its exterior square has rank
at least three, which is impossible. If $\operatorname{rank}\pi_i=2$,
choose a basis of its image. Its determinant gives a nonzero
decomposable alternating form on $U_i$, of rank two. Factoring
through $\omega_i$ would make that form a nonzero scalar multiple
of the rank-four form $\omega_i$, again impossible. Thus
$\operatorname{rank}\pi_i\le1$ for both indices, and
$\operatorname{rank}\pi\le2$.

\paragraph{2. The pure character supplies a particular nonzero evaluation.}
For an exponential character $w=r e^\ell$, direct differentiation
of the exponential gives
\[
 D_{(v,\eta)}w
   =r e^\ell(\eta+\iota_v\ell),\qquad
 c_w^2(z\wedge z')
   =r e^\ell\bigl(\pi_\ell z\wedge\pi_\ell z'\bigr),\qquad
 \pi_\ell(v,\eta)=\eta+\iota_v\ell.
 \tag{F1R9.4}
\]
The map $\pi_\ell:U\to H^1(\mathcal O_X)$ is surjective, with rank
four, and multiplication by the nonzero scalar and exponential is
injective. The preceding paragraph therefore gives
\[
                     K\not\subset\ker c_w^2.
 \tag{F1R9.5}
\]
For the unnormalized whole object (F1R9.1), replace $\ell$ by
$\ell+B$ in (F1R9.4) and carry the generators back through
$s_B^{-1}$. Formula (F1R7.5) proves, with the actual source and
target maps, the existence of
\[
 a\in HH^2(X),\qquad
 c_{\operatorname{ch}(E)}^2(a)=0,\qquad
 c_{\operatorname{ch}(Q)}^2(a)\ne0.
 \tag{F1R9.6}
\]
The full trace--Atiyah identity (F1R7.4) then implies
$\operatorname{ev}_Q(a)\ne0$.

Let $i:Q\to E$ and $p:E\to Q$ be the actual retract maps,
$pi=\operatorname{id}_Q$. Naturality of the Hochschild
characteristic action gives
\[
 p[2]\,\operatorname{ev}_E(a)\,i=\operatorname{ev}_Q(a)\ne0.
 \tag{F1R9.7}
\]
Thus $\operatorname{ev}_E(a)$ is nonzero, whereas its full trace is
zero by the first equality in (F1R9.6). Point functors act trivially
on $HH^*(X)$ and their linearizations conjugate this natural
evaluation to itself. Hence this nonzero element lies in
$\operatorname{Ext}^2_X(E,E)^G$. The actual identity-sector
induction/forgetful isomorphism in R05 identifies the invariant
ordinary trace kernel with the trace kernel of the equivariant
lift. This proves the claimed weak-semiregularity obstruction,
without reversing a merely sufficient implication.

\paragraph{3. Exact extension through the listed functor orbit.}
The degree-two intertwiners (F1R8.6)--(F1R8.9) carry the kernel of
(F1R9.4) to
\[
           \ker\bigl(\bigwedge^2(\pi_\ell M^{-1})\bigr),
 \tag{F1R9.8}
\]
where $M$ is invertible on $HH^1$. The target map is invertible
up to a nonzero scalar, including finite-isogeny pushforward.
The quotient $\pi_\ell M^{-1}$ still has rank four. Paragraph 1
and the same actual retract argument therefore apply. For a
skyscraper sheaf, the Poincar\'e Fourier--Mukai transform is a
degree-zero line bundle; its shift and inverse transform are among
these operations. For a simple semihomogeneous bundle, the full
character identity can be recovered from its actual isogeny
presentation. Mukai's published Theorem 5.8, also stated in
Orlov's Proposition 4.5, gives $Q=f_*L$ for an isogeny
$f:Y\to X$ and a line bundle $L$. Finite flat base change gives
$f^*Q=\bigoplus_{h\in\ker f}t_h^*L$, and translations fix divisor
classes. Thus
\[
 f^*\operatorname{ch}(Q)=r e^{c_1(L)},\qquad
 f^*c_1(Q)=r c_1(L).
\]
The identity $f_*f^*=(\deg f)\operatorname{id}$ makes rational
pullback injective, proving
$\operatorname{ch}(Q)=r e^{c_1(Q)/r}$. This is precisely the
full formula (4.8) in Mukai's author translation, p.~10, used
in DEG\_L S-37.

\paragraph{4. Relation to the earlier gate.}
DEG\_L S-37 excludes split sums of shifted simple bundles with
exponential slopes in $F\theta$, by the four-slope moment matrix
and equivariant slope idempotents. R09 proves an additional
unconditional restriction: one actual nonzero exponential
character summand suffices, its slope is any rational algebraic
divisor, its complement is arbitrary, and its idempotent need not
be $G$-invariant. The exact morphism replacing those slope
idempotents is (F1R9.7), followed by the invariant Hochschild
evaluation. No persistence of a kernel through an unconstructed
cone is asserted.

\subsection*{Framework maps for R07--R09}

These statements use actual perfect objects before applying their
Grothendieck-group characters. An equivalence with kernel
$\mathcal K\in\operatorname{Perf}(X\times Y)$ has the exact map
\[
 K_0(X)\longrightarrow K_0(Y),\qquad
 x\longmapsto p_{Y!}\bigl(p_X^*x\otimes[\mathcal K]\bigr).
 \tag{F1R9.9}
\]
Grothendieck--Riemann--Roch, with both Todd classes one, realizes
this as
\[
             \Phi^H(v)=p_{Y*}(p_X^*v\operatorname{ch}(\mathcal K)).
\]
It is this formula that gives the surface Euler/rank calculation
(F1R8.10), the nonzero character transports, and the relation
between the actual kernel and (F1R8.6).

The Newton exterior operations of SGA 6 V, Appendix 7.1--7.5,
realize Adams $\psi^a$ by
$\operatorname{ch}_j(\psi^a x)=a^j\operatorname{ch}_j(x)$,
as directly rederived in HD-F1-F04. In particular the exact
balanced-character formula is
\[
 \operatorname{ch}(\psi^a[E])
      =N e^{-aB}\bigl(1-a^2q(A_1+A_2)+a^4q^2A_1A_2\bigr).
 \tag{F1R9.10}
\]
The degree-$j$ projector on the character of $K_0(X)_{\mathbf Q}$
can also be written as the explicit Newton-operation polynomial
\[
 \Pi_j=\prod_{\substack{0\le k\le4\\k\ne j}}
           \frac{\psi^2-2^k}{2^j-2^k},\qquad
 \operatorname{ch}(\Pi_jx)=\operatorname{ch}_j(x).
 \tag{F1R9.11}
\]
This formula is a rational Grothendieck-group identity. The
specific characteristic maps and retracts in (F1R7.4) and
(F1R9.7) retain the actual objects on which the Ext conclusions
are proved.

For completeness the point action used throughout is supplied
by the normalized Poincar\'e biextension $\mathcal P_X$.
With the pullback convention $T_{x,\alpha}=t_x^*(-)\otimes P_\alpha$,
its rigidified translation identity is
$t_x^*P_\beta\simeq P_\beta\otimes(\mathcal P_X)_{(x,\beta)}$.
Consequently the actual functor composition is
\[
 T_{x,\alpha}T_{y,\beta}
   \simeq T_{x+y,\alpha+\beta}
                     \otimes(\mathcal P_X)_{(x,\beta)}.
 \tag{F1R9.12}
\]
The factor on the right is a one-dimensional constant vector
space, and the biextension identities give its associativity
constraints. A specified coherent finite action retains the
corresponding multiplication choices and their associativity;
none is replaced by a projective action without a lift.
Pullback along $u$ obeys the exact biextension equality
$(u\times1)^*\mathcal P_Y\simeq(1\times\widehat u)^*\mathcal P_X$,
which is the line-bundle identity behind (F1R8.7).
Conjugation by a line tensor gives (F1R8.8), and conjugation by
an actual Fourier--Mukai kernel gives Orlov's graph map
(F1R8.5). These are the SGA 7 I VII--VIII morphisms recorded
and checked in HD-F1-F03.

Finally, the passage from these operations to a nonzero trace
kernel is the published Buchweitz--Flenner evaluation/trace
diagram, with invariant Serre transposes, followed by Perry's
published finite-action Hochschild decomposition. R07 uses its
algebra structure and degree-zero trace, R08 uses its actual
degree-one point-group differential, and R09 uses its
degree-two natural evaluation on the specified retract.
